\subsection*{一步一步建立直观}

\paragraph{1. 换尺子不等于移动对象}

绿色向量代表空间对象，坐标是用选定基向量给它写出的配方。换基只改变配方的记录方式，不应凭空移动箭头。右图选取实对称矩阵，虚线提供一组特别方便的基方向，作为一般换基的一个具体例子。



\paragraph{2. 为什么某些坐标更容易计算}

当绿色箭头沿第一条特征方向时，矩阵作用只把它放大，不产生另一方向的分量。改用两条特征方向为基，每个坐标就各自乘上一个倍率。原本混合两个分量的计算，因此变成独立的缩放。



\paragraph{3. 按翻译顺序推导矩阵表达}

若 $P$ 的列是新基在旧坐标中的表达，新坐标先乘 $P$ 才能交给旧矩阵 $A$，输出再乘 $P^{-1}$ 译回。因此新矩阵是 $P^{-1}AP$。顺序来自坐标转换的职责；不是因为三块矩阵可以随意交换。


